% proflycee-tools-aleatoire.tex
% Copyright 2023-2026 Cédric Pierquet
% Released under the LaTeX Project Public License v1.3c or later, see http://www.latex-project.org/lppl.txt

%===trad fr/en ok

%%-----Ensemble, version simplifiée
\newcommand\pflstrut{\vphantom{(Q)^{2}}}
\newcommand\pflensemble[2][]{%
  \left\lbrace\mathpunct{}#1#2\mathpunct{}\right\rbrace
}

%%------EnsembleList
\defKV[kvensemble]{%
  Sep=\def\PLenssep{#1},sep=\def\PLenssep{#1},%
  Option=\def\PLensopt{#1},option=\def\PLensopt{#1}
}

\setKVdefault[kvensemble]{%
  Sep={;},sep={;},%
  Mathpunct=true,mathpunct=true,%
  Option={},option={}
}

\newcommand\EcritureEnsemble[2][]{%
  \useKVdefault[kvensemble]%
  \setKV[kvensemble]{#1}%
  \setKVbooltruedefaultmulti[kvensemble]{mathpunct}{Mathpunct}%
  \ifboolKV[kvensemble]{Mathpunct}%
    {\left\lbrace \PLensopt{} \mathpunct{} \StrSubstitute{#2}{/}{\mathpunct{}{\PLenssep}\mathpunct{}} \mathpunct{} \right\rbrace}%
    {\left\lbrace \PLensopt{} \StrSubstitute{#2}{/}{{\PLenssep}} \right\rbrace}%
}
\NewCommandCopy\pflecritens\EcritureEnsemble
\NewCommandCopy\pflwriteset\EcritureEnsemble

%%------TRINOMEALEA
\setKVdefault[trinomealea]{%
  Alea=false,alea=false,%
  Anegatif=true,aneg=true
}

\NewDocumentCommand\EcritureTrinome{ O{} m m m }{%
  \useKVdefault[trinomealea]%
  \setKV[trinomealea]{#1}%
  \setKVboolfalsedefaultmulti[trinomealea]{alea}{Alea}%
  \setKVbooltruedefaultmulti[trinomealea]{aneg}{Anegatif}%
  \ifboolKV[trinomealea]{Alea}%
  {%
    \xdef\Coeffa{\fpeval{randint(#2)}}\xdef\Coeffb{\fpeval{randint(#3)}}\xdef\Coeffc{\fpeval{randint(#4)}}%
    \xdef\AleaSigneA{\fpeval{randint(1,2)}}%
    \ifboolKV[trinomealea]{Anegatif}%
      {\xintifboolexpr{\AleaSigneA == 1}{\def\TriSigneA{}}{\def\TriSigneA{-}}}%
      {\def\TriSigneA{}}%forcément a>0
  }%
  {%
    \xdef\Coeffa{#2}\xdef\Coeffb{#3}\xdef\Coeffc{#4}%
    \def\TriSigneA{}%
  }%
  \def\TriPartieA{\num{\TriSigneA\Coeffa} x^2}%
  \xintifboolexpr{\Coeffa == 1}%
    {\def\TriPartieA{x^2}}{}%
  \xintifboolexpr{\Coeffa == -1}%
    {\def\TriPartieA{-x^2}}{}%
  \def\TriPartieB{\num[print-implicit-plus]{\Coeffb}x}%
  \xintifboolexpr{\Coeffb == 0}%
    {\def\TriPartieB{}}{}%
  \xintifboolexpr{\Coeffb == 1}%
    {\def\TriPartieB{+x}}{}%
  \xintifboolexpr{\Coeffb == -1}%
    {\def\TriPartieB{-x}}{}%
  \xintifboolexpr{\Coeffc == 0}%
    {\def\TriPartieC{}}{\def\TriPartieC{\num[print-implicit-plus]{\Coeffc}}}%
  %formatage
  \ensuremath{\TriPartieA\TriPartieB\TriPartieC}%
}
\NewCommandCopy\pflecrittrinom\EcritureTrinome
\NewCommandCopy\pfltrinomial\EcritureTrinome

%%------ALEATOIRES
\ExplSyntaxOn
\NewDocumentCommand\TriListeCroiss{m}
{
  \clist_sort:Nn #1
  {
    \fp_compare:nNnTF {##1} > {##2}
    { \sort_return_swapped: }
    { \sort_return_same: }
  }
}
\NewDocumentCommand\TriListeDecroiss{m}
{
  \clist_sort:Nn #1
  {
    \fp_compare:nNnTF {##1} < {##2}
    { \sort_return_swapped: }
    { \sort_return_same: }
  }
}
\ExplSyntaxOff

\newcommand\NbAlea[4][0]{%entier ou décimal
  \xintifboolexpr{#1 == 0}%
    {%
      \xdef#4{\fpeval{randint(#2,#3)}}%
    }%
    {%
      \xintifboolexpr{(#3-(#2)) == 1}%
        {%
          \xdef#4{\fpeval{round(#2+rand(),#1)}}%
        }%
        {%
          \xdef#4{\fpeval{round(randint(#2,#3)+rand(),#1)}}%
        }%
    }%
}
\NewCommandCopy\pflnbalea\NbAlea
\NewCommandCopy\pflrandnum\NbAlea

\newcommand\VarNbAlea[2]{%
  \xdef#1{\fpeval{#2}}%
}
\NewCommandCopy\pflvarnbalea\VarNbAlea
\NewCommandCopy\pflvarrandnum\VarNbAlea

\defKV[TirageAleatoireEntiersEntre]{%
  ValMin=\def\TAEEmin{#1},min value=\def\TAEEmin{#1},%
  ValMax=\def\TAEEmax{#1},max value=\def\TAEEmax{#1},%
  NbVal=\def\TAEEnb{#1},nb values=\def\TAEEnb{#1},%
  Sep=\def\TAEEsep{#1},sep=\def\TAEEsep{#1},%
  Tri=\def\TAEEtri{#1},sort=\def\TAEEtri{#1},%
}

\setKVdefault[TirageAleatoireEntiersEntre]{%
  ValMin=1,min value=1,%
  ValMax=50,max value=50,%
  NbVal=6,nb values=6,%
  Sep={,},sep={,},%
  Tri=non,sort=no,%
  Repetition=false,repeat=false
}

\NewList{TirageEntiersAleatoires}
\NewDocumentCommand\TirageAleatoireEntiers{ O{} m }{%1=options,2=nomliste triée ou non
  \useKVdefault[TirageAleatoireEntiersEntre]%
  \setKV[TirageAleatoireEntiersEntre]{#1}%
  %multilng booleans
  \setKVboolfalsedefaultmulti[TirageAleatoireEntiersEntre]{repeat}{Repetition}%
  %back to normal
  \ifboolKV[TirageAleatoireEntiersEntre]{Repetition}%suivant la possibilité de répéter
    {%répétitions autorisées
      %création de la liste (le 1er puis les autres)
      \xdef#2{\fpeval{randint(\TAEEmin,\TAEEmax)}}%
      \foreach \i in {2,3,...,\TAEEnb}%
        {\xdef#2{#2,\fpeval{randint(\TAEEmin,\TAEEmax)}}}%
    }%
    {%pas de répétitions
      %melange des nombres
      \ClearList{TirageEntiersAleatoires}%on vide la liste aléatoire
      \xintFor* ##1 in {\xintSeq{\TAEEmin}{\TAEEmax}}%
        \do{\InsertRandomItem{TirageEntiersAleatoires}{##1}}%
      %création de la liste (le 1er puis les autres)
      \xdef#2{\TirageEntiersAleatoires[0]}%
      \foreach \i in {1,2,...,{\inteval{\TAEEnb-1}}}%
        {\xdef#2{#2,\TirageEntiersAleatoires[\i]}}%
    }%
  %on trie si besoin [fr]/[en]
  \IfSubStr{,croissant,asc,}{,\TAEEtri,}%
  %\IfStrEq{\TAEEtri}{croissant}%si tri croissant
    {\TriListeCroiss{#2}}%
    {}%
  \IfSubStr{,decroissant,des,}{,\TAEEtri,}%
  %\IfStrEq{\TAEEtri}{decroissant}%si tri décroissant
    {\TriListeDecroiss{#2}}%
    {}%
  \StrSubstitute{#2}{,}{\TAEEsep}[#2]%on change le séparateur si besoin
}
\NewCommandCopy\pfltiragealeaent\TirageAleatoireEntiers
\NewCommandCopy\pflrandintnumbers\TirageAleatoireEntiers

\endinput